% !TEX program = xelatex
% ============================================================
%  直接粘贴至 Overleaf main.tex，Compiler 选 XeLaTeX
%  已针对免费套餐编译超时问题精简：
%    - tcolorbox(skins) → mdframed（去除 TikZ 依赖）
%    - 移除 fontawesome5
%    - 修复 headheight / hyperref 警告
% ============================================================

\begin{filecontents*}{references.bib}
@article{woodham1980photometric,
  title     = {Photometric method for determining surface orientation
               from multiple images},
  author    = {Woodham, Robert J.},
  journal   = {Optical Engineering},
  volume    = {19}, number = {1}, pages = {139--144}, year = {1980},
  doi       = {10.1117/12.7972479},
  url       = {https://doi.org/10.1117/12.7972479}
}
@article{frankot1988method,
  title     = {A method for enforcing integrability in shape from
               shading algorithms},
  author    = {Frankot, Robert T. and Chellappa, Rama},
  journal   = {IEEE Trans.\ Pattern Anal.\ Mach.\ Intell.},
  volume    = {10}, number = {4}, pages = {439--451}, year = {1988},
  doi       = {10.1109/34.3910},
  url       = {https://doi.org/10.1109/34.3910}
}
@article{simchony1990direct,
  title     = {Direct analytical methods for solving {Poisson} equations
               in computer vision problems},
  author    = {Simchony, Tal and Chellappa, Rama and Shao, Min},
  journal   = {IEEE Trans.\ Pattern Anal.\ Mach.\ Intell.},
  volume    = {12}, number = {5}, pages = {435--446}, year = {1990},
  doi       = {10.1109/34.55108},
  url       = {https://doi.org/10.1109/34.55108}
}
@inproceedings{johnson2009retrographic,
  title     = {Retrographic sensing for the measurement of surface
               texture and shape},
  author    = {Johnson, Micah K. and Adelson, Edward H.},
  booktitle = {CVPR}, pages = {1070--1077}, year = {2009},
  doi       = {10.1109/CVPR.2009.5206534},
  url       = {https://doi.org/10.1109/CVPR.2009.5206534}
}
@article{yuan2017gelsight,
  title     = {{GelSight}: High-resolution robot tactile sensors for
               estimating geometry and force},
  author    = {Yuan, Wenzhen and Dong, Siyuan and Adelson, Edward H.},
  journal   = {Sensors},
  volume    = {17}, number = {12}, pages = {2762}, year = {2017},
  doi       = {10.3390/s17122762},
  url       = {https://doi.org/10.3390/s17122762}
}
@article{lambeta2020digit,
  title     = {{DIGIT}: A novel design for a low-cost compact
               high-resolution tactile sensor},
  author    = {Lambeta, Mike and Chou, Po-Wei and Tian, Stephen and
               Yang, Brian and others},
  journal   = {IEEE Robot.\ Autom.\ Lett.},
  volume    = {5}, number = {3}, pages = {3838--3845}, year = {2020},
  doi       = {10.1109/LRA.2020.3009864},
  url       = {https://doi.org/10.1109/LRA.2020.3009864}
}
@article{queeau2018normal,
  title     = {Normal integration: {A} survey},
  author    = {Qu{\'e}au, Yvain and Durou, Jean-Denis and
               Aujol, Jean-Fran\c{c}ois},
  journal   = {J.\ Math.\ Imaging Vis.},
  volume    = {60}, number = {4}, pages = {576--593}, year = {2018},
  doi       = {10.1007/s10851-018-0819-7},
  url       = {https://doi.org/10.1007/s10851-018-0819-7}
}
@inproceedings{harker2008least,
  title     = {Least squares surface reconstruction from measured
               gradient fields},
  author    = {Harker, Matthew and O'Leary, Paul},
  booktitle = {CVPR}, pages = {1--7}, year = {2008},
  doi       = {10.1109/CVPR.2008.4587414},
  url       = {https://doi.org/10.1109/CVPR.2008.4587414}
}
@inproceedings{wang2021gelsightwedge,
  title     = {{GelSight Wedge}: Measuring high-resolution {3D} contact
               geometry with a compact robot finger},
  author    = {Wang, Shaoxiong and Lambeta, Mike and Chou, Po-Wei
               and Calandra, Roberto},
  booktitle = {ICRA}, pages = {6468--6475}, year = {2021},
  doi       = {10.1109/ICRA48506.2021.9561781},
  url       = {https://doi.org/10.1109/ICRA48506.2021.9561781}
}
@inproceedings{dong2017improved,
  title     = {Improved {GelSight} tactile sensor for measuring
               geometry and slip},
  author    = {Dong, Siyuan and Yuan, Wenzhen and Adelson, Edward H.},
  booktitle = {IROS}, pages = {137--144}, year = {2017},
  doi       = {10.1109/IROS.2017.8206534},
  url       = {https://doi.org/10.1109/IROS.2017.8206534}
}
\end{filecontents*}

% ============================================================
\documentclass[12pt, a4paper]{ctexart}

% ---- 页面 ----
\usepackage[top=2.8cm, bottom=2.8cm, left=3.0cm, right=3.0cm,
            headheight=14pt]{geometry}   % headheight 一并修复

% ---- 数学 ----
\usepackage{amsmath, amssymb}

% ---- 颜色 ----
\usepackage{xcolor}
\definecolor{codebg}   {RGB}{248,248,248}
\definecolor{codeframe}{RGB}{200,200,200}
\definecolor{kwcolor}  {RGB}{  0,  0,180}
\definecolor{cmcolor}  {RGB}{110,110,110}
\definecolor{strcolor} {RGB}{  0,130,  0}

% ---- 代码（listings，无需 minted 或额外编译步骤）----
\usepackage{listings}
\lstset{
  language         = Python,
  basicstyle       = \ttfamily\small,
  keywordstyle     = \color{kwcolor}\bfseries,
  commentstyle     = \color{cmcolor}\itshape,
  stringstyle      = \color{strcolor},
  showstringspaces = false,
  breaklines       = true,
  breakatwhitespace= true,
  frame            = single,
  framesep         = 4pt,
  rulecolor        = \color{codeframe},
  backgroundcolor  = \color{codebg},
  xleftmargin      = 1.6em,
  xrightmargin     = 0.4em,
  numbers          = left,
  numberstyle      = \tiny\color{gray},
  stepnumber       = 1,
  numbersep        = 5pt,
  captionpos       = b,
  aboveskip        = 0.6em,
  belowskip        = 0.3em,
  tabsize          = 4,
}

% ---- 彩色提示框（mdframed，轻量替代 tcolorbox+skins）----
\usepackage[framemethod=default]{mdframed}

\mdfdefinestyle{bugstyle}{
  backgroundcolor = red!4!white,
  linecolor       = red!55!black,
  linewidth       = 1.2pt,
  roundcorner     = 3pt,
  innerleftmargin = 8pt, innerrightmargin = 8pt,
  innertopmargin  = 6pt, innerbottommargin = 6pt,
  skipabove       = 6pt, skipbelow = 4pt,
}
\mdfdefinestyle{fixstyle}{
  backgroundcolor = green!4!white,
  linecolor       = green!50!black,
  linewidth       = 1.2pt,
  roundcorner     = 3pt,
  innerleftmargin = 8pt, innerrightmargin = 8pt,
  innertopmargin  = 6pt, innerbottommargin = 6pt,
  skipabove       = 6pt, skipbelow = 4pt,
}

\newenvironment{bugbox}{%
  \begin{mdframed}[style=bugstyle]%
  \noindent{\bfseries\color{red!70!black}[问题 / Bug]\quad}%
}{%
  \end{mdframed}%
}
\newenvironment{fixbox}{%
  \begin{mdframed}[style=fixstyle]%
  \noindent{\bfseries\color{green!60!black}[修复方案]\quad}%
}{%
  \end{mdframed}%
}

% ---- 表格 ----
\usepackage{booktabs, tabularx, multirow}
\renewcommand{\arraystretch}{1.25}

% ---- 超链接（禁止在 PDF 书签中出现数学符号，消除警告）----
\usepackage{hyperref}
\hypersetup{
  colorlinks  = true,
  linkcolor   = {blue!70!black},
  citecolor   = {blue!70!black},
  urlcolor    = {blue!70!black},
  pdftitle    = {GelSight 光度立体重建算法精度问题分析},
  pdfauthor   = {},
}

% ---- 参考文献 ----
\usepackage[numbers,sort&compress]{natbib}

% ---- 页眉页脚 ----
\usepackage{fancyhdr}
\pagestyle{fancy}
\fancyhf{}
\rhead{\small\itshape GelSight 光度立体算法精度分析}
\lhead{\small\itshape 内部技术报告}
\cfoot{\thepage}
\renewcommand{\headrulewidth}{0.4pt}

% ---- 列表 ----
\usepackage{enumitem}
\setlist[enumerate]{leftmargin=*, label=\arabic*.}
\setlist[itemize]{leftmargin=*}

% ---- 行距 ----
\usepackage{setspace}
\onehalfspacing

% ============================================================
\begin{document}
% ============================================================

% -------- 封面 --------
\begin{titlepage}
  \centering
  \vspace*{1.8cm}
  {\LARGE\bfseries GelSight 光度立体重建算法\\[0.5em]精度问题分析报告\par}
  \vspace{0.5em}
  {\large\itshape Accuracy Issue Analysis of the GelSight\\[0.2em]
   Photometric Stereo Reconstruction Pipeline\par}
  \vspace{1.8em}
  \rule{0.55\linewidth}{0.5pt}\\[1em]
  {\normalsize
    基于源码 \texttt{gelsight\_ps} 的逐模块审查\\[0.4em]
    覆盖模块：\texttt{geometry.py}、\texttt{ball\_calibrate.py}、\\[0.2em]
    \texttt{poisson.py}、\texttt{reconstruct.py}、\texttt{image\_loader.py}
  \par}
  \vspace{1.8em}
  {\normalsize\today}
  \vfill
  \begin{mdframed}[
    backgroundcolor = gray!7,
    linecolor       = gray!45,
    linewidth       = 0.6pt,
    roundcorner     = 4pt,
    innerleftmargin = 12pt, innerrightmargin = 12pt,
    innertopmargin  = 10pt, innerbottommargin = 10pt,
  ]
    \small\textbf{摘要：}本报告对 \texttt{gelsight\_ps} 软件包的光度立体标定与
    重建流程进行全面精度审查，识别出 \textbf{6~类、12~个具体问题}，包括
    3~个可立即修复的代码逻辑缺陷及 9~个需算法层面改进的问题。各问题均附有
    定位代码、误差分析、相关文献（含 DOI 链接）及修复建议，并按影响程度提供
    修复优先级排序。
  \end{mdframed}
\end{titlepage}

\tableofcontents
\newpage

% ============================================================
\section{引言}
\label{sec:intro}
% ============================================================

GelSight 系列触觉传感器~\cite{johnson2009retrographic,yuan2017gelsight}
通过观察弹性凝胶层在接触时的光度变化，配合光度立体（Photometric Stereo, PS）
算法~\cite{woodham1980photometric} 重建接触面的三维形貌。其标定流程的核心是：
用已知几何形状（钢球）压入传感器，建立 RGB 响应到表面法向的逐像素映射（LUT）；
推理时通过 LUT 预测法向，再经 Frankot-Chellappa 频域积分~\cite{frankot1988method}
恢复深度图。

经逐行源码审查，当前实现在以下六个方向存在精度问题。

\begin{table}[ht]
  \centering
  \caption{问题分类概览}
  \label{tab:overview}
  \begin{tabular}{clcc}
    \toprule
    \textbf{章节} & \textbf{问题类别} & \textbf{主要模块} & \textbf{数量} \\
    \midrule
    \S\ref{sec:gt_bias}     & 地面真值法向量系统性偏差   & \texttt{geometry.py}        & 2 \\
    \S\ref{sec:regression}  & 回归模型内在局限           & \texttt{ball\_calibrate.py} & 4 \\
    \S\ref{sec:integration} & 法向积分误差               & \texttt{poisson.py}         & 2 \\
    \S\ref{sec:bugs}        & 代码逻辑缺陷（可立即修复） & \texttt{reconstruct.py}     & 2 \\
    \S\ref{sec:eval}        & 评估指标非标准             & \texttt{ball\_calibrate.py} & 1 \\
    \S\ref{sec:preproc}     & 预处理缺陷                 & \texttt{image\_loader.py}   & 1 \\
    \midrule
    & \textbf{合计} & & \textbf{12} \\
    \bottomrule
  \end{tabular}
\end{table}

% ============================================================
\section{地面真值法向量的系统性偏差}
\label{sec:gt_bias}
% ============================================================

\subsection{\texttt{nz\_floor} 直接篡改训练标签（高优先级）}
\label{subsec:nz_floor}

\subsubsection*{问题定位}

\texttt{geometry.py}，函数 \texttt{sphere\_normals\_for\_contact()}：

\begin{lstlisting}[caption={\texttt{nz\_floor} 对训练标签的错误应用}]
nz_raw = z[inside] / R
# 本意是防止推理阶段梯度奇异，但错误地修改了训练标签
nz[inside] = np.maximum(nz_raw, float(nz_floor))  # nz_floor=0.12
n    = np.stack([nx, ny, nz], axis=-1)
norm = np.linalg.norm(n, axis=-1, keepdims=True) + 1e-9
n    = n / norm   # 归一化后 nz 不再是球面真值
\end{lstlisting}

\begin{bugbox}
当接触区边缘像素的 $n_z^{raw} < 0.12$ 时，代码将其强制提升至 $0.12$，
随后对向量 $(n_x,\, n_y,\, 0.12)$ 归一化，得到：
\[
  \hat{\mathbf{n}}_{label}
  = \frac{(n_x,\; n_y,\; 0.12)}{\|(n_x,\; n_y,\; 0.12)\|}
  \;\neq\; \hat{\mathbf{n}}_{true}
\]
回归器以被篡改的 $\hat{\mathbf{n}}_{label}$ 为目标训练，导致高倾斜角区域
（$\theta > \arccos(0.12) \approx 83^\circ$）法向预测出现\textbf{系统性偏差}。
\end{bugbox}

\subsubsection*{理论依据}

Woodham~\cite{woodham1980photometric} 在原始 PS 论文中指出：近掠射角样本
不确定度高，应\textbf{从训练集中剔除}，而非修正。\texttt{nz\_floor} 作为梯度
分母截断（防止 $p = n_x/n_z$ 奇异化），应仅在推理阶段 \texttt{reconstruct.py}
中使用。

\begin{fixbox}
扩大 \texttt{rim\_shrink\_pix} 使进入训练的像素均满足
$n_z^{raw} \geq \texttt{nz\_floor}$，不对法向向量本身做任何修改：

\begin{lstlisting}
# geometry.py -- 修复：排除低 nz 像素，而非修改其标签
valid_nz     = (z[inside] / R) >= float(nz_floor)
inside_valid = inside.copy()
inside_valid[inside] = valid_nz        # 只保留合格像素
mask             = inside_valid.astype(np.uint8)
nz[inside_valid] = z[inside][valid_nz] / R  # 原始球面法向，不做 floor
\end{lstlisting}
\end{fixbox}

% ----------------------------------------
\subsection{Hertz 接触变形导致球面几何模型失效}
\label{subsec:hertz}

\texttt{sphere\_normals\_for\_contact()} 假设接触区法向严格等于刚性球面法向。
然而 GelSight 所用硅胶（Shore~00-20 至 00-40 级）是弹性软材料：根据 Hertz
弹性接触理论，球体压入时接触边缘凝胶发生非线性弯折变形，真实法向与刚性球面
法向可差 $5°$--$15°$，Yuan~et~al.~\cite{yuan2017gelsight} 及
Dong~et~al.~\cite{dong2017improved} 对此均有定性讨论。

\textbf{改进方向：}
\begin{enumerate}
  \item \textbf{多形状交叉验证}：用平面、圆柱、球面等标准件交叉验证，
        当预测深度与几何真值一致时说明模型已收敛至物理正确状态。
  \item \textbf{数据驱动弹性修正}：通过对已知平面的残差法向场统计建模，
        得到与接触深度相关的修正量，叠加至 $\hat{\mathbf{n}}_{sphere}$。
\end{enumerate}

% ============================================================
\section{回归模型的内在局限}
\label{sec:regression}
% ============================================================

\subsection{线性（二阶多项式）模型无法描述复杂光学效应}
\label{subsec:linear_model}

\texttt{build\_feature\_extractor()} 构造 11 维多项式特征，对应线性回归：
\[
  \boldsymbol{\phi}(\mathbf{c})
  = \bigl[r',\; g',\; b',\; I,\; 1,\; R^2,\; G^2,\; B^2,\; RG,\; RB,\; GB\bigr],
  \qquad
  \hat{\mathbf{n}} = \mathbf{W}\boldsymbol{\phi}
\]

\begin{table}[ht]
  \centering
  \caption{GelSight 光学通道中的非线性效应}
  \label{tab:nonlinear}
  \begin{tabular}{lcc}
    \toprule
    \textbf{效应}                      & \textbf{线性可描述？} & \textbf{影响区域} \\
    \midrule
    多 LED 距离衰减（$1/d^2$ 律）      & 否        & 全局         \\
    凝胶--空气界面菲涅尔反射            & 否        & 大倾斜角区域 \\
    硅胶内次表面散射                    & 否        & 接触核心区   \\
    多 LED 混合色散差异                 & 部分      & 全局         \\
    \bottomrule
  \end{tabular}
\end{table}

Lambeta~et~al.~\cite{lambeta2020digit} 在 DIGIT 传感器中用 CNN 替代逐像素
线性回归，在高曲率区域法向误差显著降低。短期改进是扩展特征空间：

\begin{lstlisting}[caption={短期改进：追加三阶单项式}]
feature_names += ["R^3", "G^3", "B^3"]
feats += [R * R * R, G * G * G, B * B * B]
\end{lstlisting}

% ----------------------------------------
\subsection{8{\texttimes}8 分块网格过粗，无法捕捉像素级光照变化}
\label{subsec:block_grid}

\begin{lstlisting}[caption={硬编码的 8x8 分块（ball\_calibrate.py 第 1066 行）}]
K  = 8
bh = (H + K - 1) // K   # 对 480x640 图像约为 60x80 像素/块
bw = (W + K - 1) // K
\end{lstlisting}

同一块内像素到 LED 的方向角最大可相差 $20°$--$30°$，块内的 RGB→法向映射并非
同一个线性变换。Wang~et~al.~\cite{wang2021gelsightwedge} 的 GelSight~Wedge
使用逐像素 LUT，正是对此局限的直接回应。

\begin{fixbox}
将 $K$ 提高至 16 或 32；对覆盖充足像素直接逐像素岭回归：

\begin{lstlisting}
K = 16   # 块大小降至约 15x20 像素

for y, x in zip(*np.where(Ns >= min_spp)):
    W_field[y, x] = np.linalg.solve(
        XtX[y, x] + lam * np.eye(FEAT_DIM),
        XtY[y, x]
    )
# 覆盖不足的像素用邻域加权（RBF）插值填充
\end{lstlisting}
\end{fixbox}

% ----------------------------------------
\subsection{高斯融合邻域搜索范围不足}
\label{subsec:gaussian_blend}

\begin{lstlisting}[caption={sigma 与搜索范围不匹配（ball\_calibrate.py 第 1113 行）}]
sigma_x = bw * 0.75   # 有效半径约 1.5 个块宽
sigma_y = bh * 0.75
# 搜索范围仅 +-1 块（约 +-1.33 sigma），应至少 +-2 sigma
for by in range(max(0, int(by_f) - 1), min(K, int(by_f) + 2)):
    for bx in range(max(0, int(bx_f) - 1), min(K, int(bx_f) + 2)):
\end{lstlisting}

图像角点处只有 1--2 个有效邻居，$w_{sum}$ 极小，导致：
\[
  \mathbf{W}_{acc}
  = \mathbf{W}_{global}
    + \frac{\mathbf{W}_{acc} - \mathbf{W}_{global}}{w_{sum}}
  \;\xrightarrow{w_{sum}\to 0}\; \text{残差被无限放大}
\]

\begin{fixbox}
搜索范围扩展至 $\pm 2$ 块，增加退化保护：

\begin{lstlisting}
radius = 2
for by in range(max(0, int(by_f) - radius),
                min(K, int(by_f) + radius + 1)):
    for bx in range(max(0, int(bx_f) - radius),
                    min(K, int(bx_f) + radius + 1)):
        ...
if wsum > 1e-6:
    W_field[y, x] = W_global + (W_acc - W_global) / wsum
else:
    W_field[y, x] = W_global   # 退化为全局模型
\end{lstlisting}
\end{fixbox}

% ----------------------------------------
\subsection{样本积累的像素级 Python 循环效率极低}
\label{subsec:loop}

\texttt{ball\_calibrate.py} 第~1009--1026~行通过 Python 循环逐像素累积
$\mathbf{X}^T\mathbf{X}$ 和 $\mathbf{X}^T\mathbf{Y}$，每帧约 3\,000--10\,000
像素，全部在 Python 层迭代。此问题\textbf{不影响精度}，但速度可提升 10--50 倍：

\begin{lstlisting}[caption={向量化方案（取代像素级 for 循环）}]
ys, xs = np.where(core > 0)
v   = X[ys, xs]          # [N, F]
yv  = Y[ys, xs]          # [N, out_dim]
w   = weights[ys, xs]    # [N]
np.add.at(XtX, (ys, xs),
          np.einsum('ni,nj->nij', v * w[:, None], v))
np.add.at(XtY, (ys, xs),
          np.einsum('ni,nj->nij', v * w[:, None], yv))
Ns[ys, xs] += (w > 0).astype(np.int32)
\end{lstlisting}

% ============================================================
\section{法向积分误差}
\label{sec:integration}
% ============================================================

\subsection{Frankot-Chellappa 的周期性假设引发低频振铃}
\label{subsec:fc_ringing}

\texttt{poisson.py} 实现的 Frankot-Chellappa 积分~\cite{frankot1988method}
通过 2D FFT 求解 Poisson 方程：
\[
  Z(\omega_x,\omega_y)
  = \frac{-j\omega_x P - j\omega_y Q}{\omega_x^2+\omega_y^2}
\]
FFT 隐含地假设信号在边界处首尾周期连接，而深度图并不满足此条件，导致低频
振铃伪影（Gibbs 效应）。Qu\'eau~et~al.~\cite{queeau2018normal} 的综述推荐基于
DCT-II 的 Poisson 求解器~\cite{simchony1990direct}，其满足 Neumann 零通量
边界条件从根本上消除振铃；Harker \& O'Leary~\cite{harker2008least}
进一步提出最小二乘最优积分方法。

\begin{fixbox}[frametitle={}]
替换 \texttt{poisson.py} 中的积分核心为 DCT-II 方案：

\begin{lstlisting}[caption={基于 DCT-II 的 Poisson 求解器（满足 Neumann BC）}]
from scipy.fft import dctn, idctn

def poisson_dct(p: np.ndarray, q: np.ndarray) -> np.ndarray:
    """Reference: Simchony et al., IEEE TPAMI, 1990."""
    H, W = p.shape
    div          = np.zeros((H, W), dtype=np.float64)
    div[:, 1:]  += p[:, :-1];  div[:, :-1]  -= p[:, :-1]
    div[1:, :]  += q[:-1, :];  div[:-1, :]  -= q[:-1, :]
    D            = dctn(div, norm='ortho')
    cy  = np.cos(np.pi * np.arange(H) / H)[:, None]
    cx  = np.cos(np.pi * np.arange(W) / W)[None, :]
    den = 2.0 * (cy - 1.0) + 2.0 * (cx - 1.0);  den[0,0] = 1.0
    D  /= den;  D[0, 0] = 0.0
    return idctn(D, norm='ortho').astype(np.float32)
\end{lstlisting}
\end{fixbox}

% ----------------------------------------
\subsection{低 \texorpdfstring{$n_z$}{nz} 处梯度放大污染全局积分}
\label{subsec:nz_amplification}

\begin{lstlisting}[caption={nz=0.1 时梯度被放大 10 倍（reconstruct.py）}]
p_grad = n[..., 0] / np.maximum(n[..., 2], 0.1)
q_grad = n[..., 1] / np.maximum(n[..., 2], 0.1)
\end{lstlisting}

当 $n_z = 0.1$（约 $84°$ 倾角）时梯度放大 10 倍。由于 FFT 是全局算子，
少量噪声大的边缘梯度会扩散至整幅深度图。
Harker \& O'Leary~\cite{harker2008least} 与
Qu\'eau~et~al.~\cite{queeau2018normal} 推荐以 $c = n_z^2$ 为置信权重降权：

\begin{fixbox}
\begin{lstlisting}
conf    = np.clip(n[..., 2], 0.0, 1.0) ** 2   # 置信度权重
nz_safe = np.maximum(n[..., 2], 0.25)          # 更保守截断（约 76 度）
p_grad  = (n[..., 0] / nz_safe) * conf
q_grad  = (n[..., 1] / nz_safe) * conf
\end{lstlisting}
\end{fixbox}

% ============================================================
\section{代码逻辑缺陷}
\label{sec:bugs}
% ============================================================

\subsection{\texttt{image\_dir} 模式下掩膜被施加两次}
\label{subsec:double_mask}

\begin{lstlisting}[caption={接触掩膜被连续施加两次（reconstruct.py）}]
# 第一次：对绝对深度施加掩膜
if contact_mask_enable and ref_lin is not None:
    mask = _residual_mask(...)
    if mask is not None:
        depth = depth * (mask / 255.0)     # depth 已被掩膜

d_diff = depth - ref_depth                 # depth 已不完整

# 第二次：对 d_diff 再次施加掩膜
if contact_mask_enable and ref_lin is not None:
    mask = _residual_mask(...)
    if mask is not None:
        d_diff = d_diff * (mask / 255.0)   # 等效二次掩膜
\end{lstlisting}

\begin{bugbox}
实际计算结果为：
$d\_diff = (depth_{orig} \cdot M - ref\_depth) \cdot M$，
正确结果应为：
$d\_diff = (depth_{orig} - ref\_depth) \cdot M$。
在掩膜外缘，第一次掩膜将 $depth_{orig}$ 置零，当 $ref\_depth \neq 0$ 时
掩膜边界处会引入假深度差。
\end{bugbox}

\begin{fixbox}
移除对绝对 \texttt{depth} 的第一次掩膜，仅对 \texttt{d\_diff} 施加一次：

\begin{lstlisting}
d_diff = (depth - ref_depth).astype(np.float32)

if contact_mask_enable and ref_lin is not None:
    mask = _residual_mask(bgr_lin, ref_lin,
                          cm_w_chroma, cm_w_int,
                          cm_thr_rel, cm_min_area)
    if mask is not None:
        d_diff *= (mask.astype(np.float32) / 255.0)
\end{lstlisting}
\end{fixbox}

% ----------------------------------------
\subsection{\texttt{out\_dim=2} 路径的 \texttt{sqrt} 产生静默 \texttt{NaN}}
\label{subsec:sqrt_nan}

\begin{lstlisting}[caption={clip 在 sqrt 之后无法清除 NaN（reconstruct.py）}]
nz = np.sqrt(1.0 - nx * nx - ny * ny)  # nx^2+ny^2>1 时 -> NaN
nz = np.clip(nz, 0.0, 1.0)             # clip(NaN) 仍然是 NaN!
\end{lstlisting}

\begin{bugbox}
当回归器外推至 $n_x^2 + n_y^2 > 1$ 时，\texttt{np.sqrt} 返回 \texttt{NaN}，
\textbf{\texttt{np.clip} 不清除 \texttt{NaN}}，其沿
normalize $\to$ $p$~grad $\to$ $q$~grad $\to$ FFT 全链路传播，
最终深度图出现随机 \texttt{NaN} 块。
注：\texttt{ball\_calibrate.py} 第~1258~行评估路径已正确处理此问题，
但 \texttt{reconstruct.py} 未同步。
\end{bugbox}

\begin{fixbox}
将 \texttt{clip} 移入 \texttt{sqrt} 内部（与 \texttt{ball\_calibrate.py} 对齐）：

\begin{lstlisting}
nz = np.sqrt(np.clip(1.0 - nx * nx - ny * ny, 0.0, 1.0))
nz = np.clip(nz, 0.0, 1.0)   # 保留作防御性检查
\end{lstlisting}
\end{fixbox}

% ============================================================
\section{评估指标问题}
\label{sec:eval}
% ============================================================

\texttt{ball\_calibrate.py} 第~1275~行使用自定义指标
$\tfrac{1+\cos\theta}{2} \in [0,1]$，
在 PS 文献中不通用且对大角度误差不敏感（$90°$ 误差时值为 $0.5$，看似"尚可"）。
PS 领域自 Woodham~\cite{woodham1980photometric} 起的标准指标为
\textbf{平均角度误差（MAE）}，单位为度：
\[
  \mathrm{MAE}
  = \frac{1}{|\mathcal{M}|}
    \sum_{(u,v)\in\mathcal{M}}
    \arccos\!\Bigl(
      \hat{\mathbf{n}}_{pred}^{(u,v)} \cdot \hat{\mathbf{n}}_{gt}^{(u,v)}
    \Bigr)
\]

\begin{fixbox}
\begin{lstlisting}
err_deg    = np.degrees(np.arccos(np.clip(cos, -1.0, 1.0)))
mae        = float(err_deg[mask].mean())
median_err = float(np.median(err_deg[mask]))
for thr in [5.0, 10.0, 15.0]:
    print(f"  < {thr:.0f} deg: {(err_deg[mask]<thr).mean()*100:.1f}%")
\end{lstlisting}
\end{fixbox}

% ============================================================
\section{预处理缺陷：缺少平场校正}
\label{sec:preproc}
% ============================================================

\texttt{Preprocessor.apply()} 执行暗场减法和 gamma 线性化，但未校正 LED 照明的
空间不均匀性（晕影，vignetting）。GelSight 图像中心与边缘的照射强度比通常达
2--4 倍~\cite{yuan2017gelsight,dong2017improved}，无法被逐像素 LUT 完全吸收。

\begin{fixbox}
对平坦未受力凝胶面拍摄参考图 $\mathbf{I}_{flat}$，在暗/白场校正后追加平场除法：

\begin{lstlisting}
# Preprocessor.__init__
if cfg.get("flat_field_path"):
    self.flat_field = self._dark_white_correct(
        imread_color(cfg["flat_field_path"]))
else:
    self.flat_field = None

# Preprocessor.apply（暗/白场校正之后）
if self.flat_field is not None:
    bgr_lin = np.clip(bgr_lin / (self.flat_field + 1e-4), 0.0, None)
\end{lstlisting}
\end{fixbox}

% ============================================================
\section{问题汇总与修复优先级}
\label{sec:summary}
% ============================================================

\begin{table}[ht]
  \centering
  \caption{全部 12 个问题汇总}
  \label{tab:summary}
  \small
  \begin{tabularx}{\textwidth}{clXcc}
    \toprule
    \textbf{\#} & \textbf{主要文件}
                & \textbf{问题描述}
                & \textbf{精度影响}
                & \textbf{修复难度} \\
    \midrule
    10 & \texttt{reconstruct.py}     & \texttt{sqrt} 产生静默 NaN
       & \textbf{极高}   & 极低 \\
    9  & \texttt{reconstruct.py}     & 掩膜施加两次，\texttt{d\_diff} 错误
       & 高              & 低   \\
    1  & \texttt{geometry.py}        & \texttt{nz\_floor} 污染训练标签
       & 高              & 低   \\
    7  & \texttt{poisson.py}         & FC 周期假设引发低频振铃
       & 高              & 中   \\
    8  & \texttt{reconstruct.py}     & 低 $n_z$ 梯度放大，污染全局积分
       & 高              & 中   \\
    5  & \texttt{ball\_calibrate.py} & 高斯融合邻域范围不足
       & 中              & 低   \\
    4  & \texttt{ball\_calibrate.py} & $K=8$ 分块过粗
       & 中高            & 中   \\
    11 & \texttt{ball\_calibrate.py} & 评估指标非标准 MAE
       & 可观测性差      & 低   \\
    12 & \texttt{image\_loader.py}   & 无平场校正
       & 中              & 中   \\
    2  & \texttt{geometry.py}        & Hertz 变形未建模
       & 中              & 高   \\
    3  & \texttt{ball\_calibrate.py} & 线性多项式模型不足
       & 中高            & 高   \\
    6  & \texttt{ball\_calibrate.py} & 像素循环低效（不影响精度）
       & 无              & 低   \\
    \bottomrule
  \end{tabularx}
\end{table}

\noindent\textbf{建议修复路线：}

\begin{enumerate}
  \item \textbf{立即可修}（不改变算法结构）\\
        \#10（\texttt{sqrt~NaN}）$\to$
        \#9（双重掩膜）$\to$
        \#11（评估指标改 MAE）$\to$
        \#6（向量化加速）

  \item \textbf{短期改进}（修改参数逻辑）\\
        \#1（\texttt{nz\_floor} 标签修复）$\to$
        \#5（扩大高斯融合邻域）

  \item \textbf{中期改进}（替换子模块）\\
        \#7（换 DCT 积分器）$\to$
        \#8（置信度加权梯度）$\to$
        \#12（添加平场校正）

  \item \textbf{长期改进}（重构算法）\\
        \#4（增大 $K$ 至逐像素回归）$\to$
        \#3（引入 MLP 非线性模型）$\to$
        \#2（Hertz 弹性修正）
\end{enumerate}

% ============================================================
\newpage
\bibliographystyle{unsrtnat}
\bibliography{references}
% ============================================================

\end{document}
